\section*{Major Revision Checklist for \textit{Bodhi VLM}}

\subsection*{How to Use This Checklist}
This checklist is organized by manuscript section and focuses only on \textbf{major-revision-level issues}. Each item is written so it can be directly used for revision planning, author rebuttal preparation, and internal paper restructuring. The goal is to move the paper from a \emph{promising framework paper} toward a \emph{top-tier journal submission} with a defensible problem definition, stronger evidence, and a tighter claim--evidence match.

\bigskip

\noindent
\textbf{Important note.} I can still prepare this checklist based on the paper I just reviewed in this conversation. However, the uploaded PDF is no longer available to reopen on my side. Therefore, the checklist below is based on the manuscript content I already analyzed, not on a fresh re-parse of the expired file. If you want, you can re-upload the PDF afterward and I can produce a second-pass version with even more line-specific comments.

\section{Title and Overall Positioning}

\subsection*{Major Revision Item 1: Narrow the paper's central claim}
\textbf{Problem.} The current title and manuscript framing suggest that the paper performs \emph{privacy budget assessment} in a broad sense. However, the actual method evaluates a \emph{feature-level discrepancy signal under an evaluator-specified perturbation model}. This is substantially narrower than estimating a true privacy budget, certifying differential privacy, or auditing end-to-end multimodal privacy guarantees.

\textbf{Why this matters.} A top-tier journal will reject papers whose title and abstract imply a stronger technical claim than the experiments actually support.

\textbf{Required revision.}
\begin{itemize}
    \item Revise the title so that it reflects \emph{feature-level budget-alignment auditing} rather than generic \emph{privacy budget assessment}.
    \item Replace broad references to ``vision and vision-language models'' with the more precise scope:
    \[
    \text{vision backbones and VLM vision encoders}.
    \]
    \item Avoid wording that implies formal privacy certification, effective $\epsilon$ estimation, or full-pipeline VLM auditing unless such claims are explicitly proven.
\end{itemize}

\textbf{Suggested title direction.}
\begin{quote}
\textit{Bodhi: Feature-Level Budget-Alignment Auditing for Vision Backbones and VLM Vision Encoders}
\end{quote}

\section{Abstract}

\subsection*{Major Revision Item 2: Rewrite the abstract to eliminate claim inflation}
\textbf{Problem.} The abstract currently reads as if the method can assess privacy budgets for vision and VLM systems in a fairly general sense. The experimental evidence does not support that broad interpretation.

\textbf{Why this matters.} The abstract is where reviewers first assess whether the paper is overclaiming.

\textbf{Required revision.}
\begin{itemize}
    \item Explicitly define the paper as an \emph{empirical feature-level auditing framework}.
    \item State that the audit is conducted \emph{under a reference perturbation model specified by the evaluator}.
    \item State clearly that the paper \textbf{does not} provide:
    \begin{itemize}
        \item formal differential privacy certification,
        \item end-to-end VLM privacy guarantees,
        \item a universal estimator of the effective privacy budget.
    \end{itemize}
    \item Replace generic statements such as ``assesses privacy budgets'' with more defensible language such as
    \[
    \text{audits alignment between declared perturbation settings and observed feature-level perturbation patterns}.
    \]
\end{itemize}

\textbf{Suggested abstract-level formalization.}
Define the paper's task as:
\[
\mathcal{A}_{\text{audit}}
=
\mathcal{A}\!\left(
\tilde{f}, f, S, \epsilon ; \mathcal{M}_{\mathrm{ref}}, \Pi, \phi
\right),
\]
where
\begin{itemize}
    \item $f$ is the original feature representation,
    \item $\tilde{f}$ is the perturbed feature representation,
    \item $S$ is the evaluator-defined sensitive subset or concept set,
    \item $\epsilon$ is the declared privacy parameter,
    \item $\mathcal{M}_{\mathrm{ref}}$ is the reference perturbation model,
    \item $\Pi$ is the feature grouping/correspondence operator,
    \item $\phi$ denotes the fitting and preprocessing protocol.
\end{itemize}

\section{Introduction}

\subsection*{Major Revision Item 3: Add a rigorous problem-definition paragraph}
\textbf{Problem.} The introduction motivates privacy auditing well, but it does not define the exact mathematical object that the paper estimates. As a result, the manuscript drifts among three notions:
\begin{enumerate}
    \item declared privacy budget $\epsilon$,
    \item perturbation scale induced by a reference mechanism,
    \item empirical discrepancy output by EMPA.
\end{enumerate}

\textbf{Why this matters.} If these three are not separated cleanly, the reader will assume the paper estimates or verifies true privacy loss.

\textbf{Required revision.}
Insert a formal problem-definition paragraph near the end of the introduction.

\textbf{Recommended LaTeX formulation.}
\[
\textbf{Input: } (f,\tilde{f},S,\epsilon,\mathcal{M}_{\mathrm{ref}}),
\qquad
\textbf{Output: } \mathrm{BAS}_\epsilon,
\]
where $\mathrm{BAS}_\epsilon$ is a \emph{budget-alignment score} measuring the discrepancy between the observed perturbed feature distribution and a reference perturbation distribution induced by $\mathcal{M}_{\mathrm{ref}}$ under declared budget $\epsilon$.

Then explicitly state:
\[
\mathrm{BAS}_\epsilon \neq \epsilon,
\qquad
\mathrm{BAS}_\epsilon \not\Rightarrow (\epsilon,\delta)\text{-DP certification}.
\]

\subsection*{Major Revision Item 4: Add a paragraph clarifying scope limitations}
\textbf{Problem.} The current introduction gives the impression that the method applies to full VLM pipelines.

\textbf{Required revision.}
Add a limitation paragraph that states:
\begin{itemize}
    \item For multimodal models, the current study evaluates \textbf{vision encoders / vision towers}, not the full language-generation pipeline.
    \item The method audits feature perturbation alignment, not semantic leakage of generated text.
    \item End-to-end multimodal privacy remains future work.
\end{itemize}

\section{Related Work}

\subsection*{Major Revision Item 5: Reorganize related work by task, not by model family}
\textbf{Problem.} The current related-work discussion risks comparing fundamentally different problem settings under one umbrella.

\textbf{Why this matters.} Reviewers will reject scope comparisons that blur distinctions between training-time DP, inference-time masking, de-identification, and post hoc auditing.

\textbf{Required revision.}
Split the related work into four clearly separated categories:
\begin{enumerate}
    \item Formal privacy mechanisms for vision and multimodal learning.
    \item Empirical privacy auditing / leakage estimation methods.
    \item Sensitive-region localization and feature attribution methods.
    \item Distribution discrepancy estimation and two-sample comparison methods.
\end{enumerate}

\textbf{Add a scope-positioning sentence such as:}
\[
\text{Bodhi is not a training-time privacy mechanism; it is a post hoc feature-level auditing framework.}
\]

\subsection*{Major Revision Item 6: Add missing baseline families in the literature discussion}
\textbf{Problem.} The paper compares mainly against generic distances such as K-L divergence or MMD, but does not sufficiently position itself against task-relevant alternatives.

\textbf{Required revision.}
Add discussion of the following baseline families:
\begin{itemize}
    \item direct perturbation-scale estimation,
    \item likelihood-based mechanism fitting,
    \item classifier two-sample testing,
    \item Wasserstein / energy distance methods,
    \item simple regression-based budget prediction from moments or summary statistics.
\end{itemize}

\section{Problem Formulation and Notation}

\subsection*{Major Revision Item 7: Separate the sensitive set from the privacy budget in notation}
\textbf{Problem.} Notation of the form $S_\epsilon$ suggests that the sensitive set depends on $\epsilon$, although in the current methodology the sensitive subset is largely evaluator-defined or detector-induced rather than budget-conditioned.

\textbf{Why this matters.} This notation creates conceptual confusion.

\textbf{Required revision.}
Use notation such as
\[
S^{(\mathrm{sens})}
\quad \text{or} \quad
S_{\mathrm{sens}}
\]
instead of $S_\epsilon$, unless the manuscript explicitly proves a budget-dependent sensitive-set construction.

\subsection*{Major Revision Item 8: Define the audit object and reference model explicitly}
\textbf{Required revision.}
Insert a dedicated definition block.

\textbf{Suggested LaTeX text.}
\begin{align}
f &\in \mathbb{R}^{d}, \\
\tilde{f} &= T_{\mathrm{obs}}(f), \\
\tilde{f}^{\,\mathrm{ref}}_{\epsilon} &= T_{\mathrm{ref},\epsilon}(f),
\end{align}
where $T_{\mathrm{obs}}$ denotes the observed perturbation process and $T_{\mathrm{ref},\epsilon}$ denotes an evaluator-specified reference perturbation mechanism at declared budget $\epsilon$.

The auditing task is then to quantify the discrepancy
\[
D_\epsilon
=
D\!\left(
\mathbb{P}(\tilde{f}\mid f,S_{\mathrm{sens}}),
\mathbb{P}(\tilde{f}^{\,\mathrm{ref}}_{\epsilon}\mid f,S_{\mathrm{sens}})
\right),
\]
or its structured approximation induced by BUA/TDA and EMPA.

\section{Method Section: BUA and TDA}

\subsection*{Major Revision Item 9: Justify the sensitive-partition heuristic}
\textbf{Problem.} BUA and TDA are currently presented as if they naturally identify sensitive features, but they are in fact heuristic grouping procedures driven by task-specific scores, thresholds, and correspondence rules.

\textbf{Why this matters.} If the partition is unreliable, all downstream EMPA outputs are unstable and hard to interpret.

\textbf{Required revision.}
Add a subsection titled something like:
\[
\text{``Why BUA/TDA is a reasonable approximation to sensitive feature localization.''}
\]

This subsection should explain:
\begin{itemize}
    \item what signal is used to define sensitivity,
    \item why that signal should correspond to privacy-relevant content,
    \item what assumptions are made,
    \item what failure cases are expected.
\end{itemize}

\subsection*{Major Revision Item 10: State the assumptions underlying cross-layer correspondence}
\textbf{Problem.} The cross-layer mapping operator is currently heuristic and insufficiently justified.

\textbf{Required revision.}
Explicitly state the approximation:
\[
\Pi_{i \rightarrow i-1}: \mathcal{F}_i \rightarrow \mathcal{F}_{i-1}
\]
is not an exact semantic correspondence operator, but a practical approximation used to propagate sensitivity across layers.

Then add an assumption block such as:
\[
\textbf{Assumption A1:}
\quad
\Pi_{i \rightarrow i-1}
\text{ preserves coarse sensitive/non-sensitive structure sufficiently for auditing purposes.}
\]

\subsection*{Major Revision Item 11: Add complexity analysis for BUA/TDA}
\textbf{Problem.} A top-tier journal expects runtime and scalability discussion.

\textbf{Required revision.}
Provide per-layer and total complexity.

For example, if the grouping plus traversal cost is dominated by pairwise operations, state a bound such as:
\[
\mathcal{O}\!\left(\sum_{i=1}^{L} n_i^2\right)
\quad \text{or} \quad
\mathcal{O}\!\left(\sum_{i=1}^{L} n_i k_i\right),
\]
depending on the actual implementation.

\section{Method Section: EMPA}

\subsection*{Major Revision Item 12: Strengthen the formal definition of EMPA}
\textbf{Problem.} The current presentation suggests that EMPA captures budget alignment, but the equations and experiments indicate that the main reported quantity may be dominated by group structure or mixture assignment rather than perturbation strength.

\textbf{Required revision.}
Give a complete hierarchical definition of EMPA.

A recommended structure is:
\begin{align}
\Theta_{\epsilon}^{\mathrm{obs}}
&=
\left(
\lambda^{\mathrm{obs}}_{\epsilon},
\mu^{\mathrm{obs}}_{\epsilon},
\Sigma^{\mathrm{obs}}_{\epsilon}
\right), \\
\Theta_{\epsilon}^{\mathrm{ref}}
&=
\left(
\lambda^{\mathrm{ref}}_{\epsilon},
\mu^{\mathrm{ref}}_{\epsilon},
\Sigma^{\mathrm{ref}}_{\epsilon}
\right),
\end{align}
and define the budget-alignment score as
\begin{equation}
\mathrm{BAS}_{\epsilon}
=
d_{\lambda}\!\left(
\lambda^{\mathrm{obs}}_{\epsilon},
\lambda^{\mathrm{ref}}_{\epsilon}
\right)
+
\alpha \,
d_{\mu}\!\left(
\mu^{\mathrm{obs}}_{\epsilon},
\mu^{\mathrm{ref}}_{\epsilon}
\right)
+
\beta \,
d_{\Sigma}\!\left(
\Sigma^{\mathrm{obs}}_{\epsilon},
\Sigma^{\mathrm{ref}}_{\epsilon}
\right),
\label{eq:bas_full}
\end{equation}
where $\alpha,\beta \ge 0$ are predefined or cross-validated weights.

\textbf{Reason.} This avoids reducing the audit to mixture weights only.

\subsection*{Major Revision Item 13: Add an identifiability subsection}
\textbf{Problem.} The empirical results suggest that EMPA bias may remain nearly constant across different $\epsilon$ values. This raises a major identifiability concern.

\textbf{Required revision.}
Add a subsection formally discussing identifiability. At minimum, the paper should ask:

\begin{quote}
Under what conditions does
\[
\mathrm{BAS}_{\epsilon_1} \neq \mathrm{BAS}_{\epsilon_2}
\]
whenever $\epsilon_1 \neq \epsilon_2$?
\end{quote}

A minimal proposition-style statement would be useful.

\textbf{Suggested mathematical direction.}
Assume a perturbation family parameterized by variance scale $\sigma^2(\epsilon)$ and a fitted family of mixture parameters $\Theta_\epsilon$. Then discuss conditions under which the map
\[
\epsilon \mapsto \Theta_\epsilon
\]
is injective, or at least rank-preserving, over the operating range.

Even if a full theorem is not possible, state clearly:
\begin{itemize}
    \item when identifiability is expected,
    \item when it may fail,
    \item what role partition quality plays.
\end{itemize}

\subsection*{Major Revision Item 14: Explain why EM is necessary}
\textbf{Problem.} A reviewer may ask why EM-based parameter analysis is needed at all instead of direct moment matching or simple scale estimation.

\textbf{Required revision.}
Add a paragraph justifying EM as modeling heterogeneous perturbation structure, for example:
\[
\mathbb{P}(\Delta f)
=
\sum_{k=1}^{K}
\lambda_k \,
\mathcal{N}(\Delta f; \mu_k, \Sigma_k),
\qquad
\Delta f = \tilde{f} - f.
\]
Then explain what information is captured by the latent mixture structure that a single-variance model would miss.

\section{Experimental Design}

\subsection*{Major Revision Item 15: Redesign the experiments around three questions}
\textbf{Problem.} The current experiments feel like demonstrations rather than decisive validation.

\textbf{Required revision.}
Reorganize the experimental section around the following three explicit research questions:

\begin{itemize}
    \item \textbf{RQ1: Identifiability.} Does the proposed score change reliably with perturbation magnitude and mechanism mismatch?
    \item \textbf{RQ2: Partition validity.} Do BUA/TDA identify privacy-relevant feature subsets more accurately than simpler alternatives?
    \item \textbf{RQ3: Practical value.} Does the audit score correlate with actual privacy or utility outcomes?
\end{itemize}

This change alone will make the paper look much more like a top-tier journal submission.

\subsection*{Major Revision Item 16: Add budget-sensitivity experiments}
\textbf{Problem.} The current results do not convincingly show that EMPA tracks budget changes.

\textbf{Required revision.}
Add controlled experiments with fixed partition and varying perturbation scale:
\[
\epsilon \in \{ \epsilon_1, \epsilon_2, \ldots, \epsilon_m \},
\qquad
\epsilon_1 < \epsilon_2 < \cdots < \epsilon_m.
\]
Then report whether
\[
\mathrm{BAS}_{\epsilon_1}, \mathrm{BAS}_{\epsilon_2}, \ldots, \mathrm{BAS}_{\epsilon_m}
\]
exhibit monotonicity, ordinal consistency, or at least statistically significant separation.

Useful summary metrics include:
\begin{align}
\rho_{\mathrm{Spearman}}
&=
\mathrm{corr}_{\mathrm{rank}}
\bigl(
\epsilon^{-1},
\mathrm{BAS}_\epsilon
\bigr), \\
\tau_{\mathrm{Kendall}}
&=
\mathrm{Kendall}
\bigl(
\epsilon^{-1},
\mathrm{BAS}_\epsilon
\bigr).
\end{align}

\subsection*{Major Revision Item 17: Add mechanism-mismatch experiments}
\textbf{Problem.} The paper currently does not show that the method distinguishes one perturbation family from another.

\textbf{Required revision.}
Add experiments of the form:
\[
T_{\mathrm{obs}} \neq T_{\mathrm{ref},\epsilon},
\]
for example:
\begin{itemize}
    \item observed Gaussian vs.\ reference Laplace,
    \item observed Laplace vs.\ reference Gaussian,
    \item observed strong perturbation vs.\ reference weak perturbation.
\end{itemize}

Then evaluate whether the audit score detects mismatch:
\[
\mathrm{BAS}\bigl(T_{\mathrm{obs}},T_{\mathrm{ref}}\bigr)
>
\mathrm{BAS}\bigl(T_{\mathrm{obs}},T_{\mathrm{obs}}\bigr).
\]

\subsection*{Major Revision Item 18: Add partition-ablation experiments}
\textbf{Problem.} It is unclear whether performance comes from BUA/TDA or simply from any nontrivial partition.

\textbf{Required revision.}
Compare against:
\begin{itemize}
    \item random partition,
    \item top-$k$ sensitivity thresholding,
    \item $k$-means on features,
    \item agglomerative clustering,
    \item saliency/attention-rollout-guided partitioning.
\end{itemize}

Then quantify partition validity using overlap metrics when a sensitive-region reference is available:
\begin{align}
\mathrm{IoU}
&=
\frac{|P \cap G|}{|P \cup G|}, \\
\mathrm{Dice}
&=
\frac{2|P \cap G|}{|P| + |G|},
\end{align}
where $P$ is the predicted sensitive subset and $G$ is the reference subset.

\subsection*{Major Revision Item 19: Add threshold-sensitivity analysis}
\textbf{Problem.} The method appears to depend strongly on quantile thresholds and grouping choices.

\textbf{Required revision.}
Run sensitivity analysis over
\[
\tau \in \{0.80, 0.85, 0.90, 0.95\}
\]
or the actual set used in the implementation, and show how audit outputs vary.

\subsection*{Major Revision Item 20: Add true task-relevant baselines}
\textbf{Problem.} Current baselines are mostly generic discrepancy measures.

\textbf{Required revision.}
Add at least the following baselines.

\paragraph{(a) Direct perturbation-scale estimation.}
Assume
\[
\Delta f = \tilde{f} - f \sim \mathcal{N}(0,\sigma^2 I)
\]
or another reference family, estimate $\hat{\sigma}$ by maximum likelihood, then map to declared budget under the assumed calibration rule.

A simple Gaussian MLE is:
\[
\hat{\sigma}^2
=
\frac{1}{nd}
\sum_{j=1}^{n}
\|\Delta f_j\|_2^2.
\]

\paragraph{(b) Moment-regression baseline.}
Construct summary statistics
\[
z =
\bigl[
\mathrm{mean}(\Delta f),
\mathrm{var}(\Delta f),
\mathrm{kurt}(\Delta f),
\mathrm{gap}_{\mathrm{sens,non\mbox{-}sens}}
\bigr]
\]
and regress to budget or budget-alignment classes.

\paragraph{(c) Two-sample classifier baseline.}
Train a classifier to distinguish observed perturbed features from reference perturbed features and use held-out AUC as a discrepancy proxy.

\paragraph{(d) Wasserstein / energy distance baseline.}
Include one or more stronger discrepancy baselines:
\begin{align}
W_1(P,Q)
&=
\inf_{\gamma \in \Gamma(P,Q)}
\int \|x-y\| \, d\gamma(x,y), \\
\mathcal{E}(P,Q)
&=
2\mathbb{E}\|X-Y\|
-
\mathbb{E}\|X-X'\|
-
\mathbb{E}\|Y-Y'\|.
\end{align}

\section{Statistical Rigor}

\subsection*{Major Revision Item 21: Convert all main results to multi-seed reporting}
\textbf{Problem.} Many key results appear to rely on one run or very limited repetitions.

\textbf{Why this matters.} With stochastic perturbation, grouping, and EM fitting, single-run results are not sufficient for a top-tier journal.

\textbf{Required revision.}
For every main table and figure, report:
\[
\text{mean} \pm \text{standard deviation}
\]
over at least $N$ seeds, where $N$ is explicitly stated.

Also add $95\%$ confidence intervals:
\[
\bar{x} \pm 1.96 \frac{s}{\sqrt{N}}.
\]

\subsection*{Major Revision Item 22: Add significance testing for key comparisons}
\textbf{Required revision.}
For paired results, report statistical significance using an appropriate test, for example:
\begin{itemize}
    \item paired $t$-test if justified,
    \item Wilcoxon signed-rank test otherwise.
\end{itemize}

If comparing multiple methods across datasets, add a correction procedure such as Holm--Bonferroni.

\section{Results Interpretation}

\subsection*{Major Revision Item 23: Address the near-constant EMPA outputs directly}
\textbf{Problem.} If EMPA or its bias term is nearly unchanged across budgets, this must be treated as a central methodological issue, not a minor observation.

\textbf{Required revision.}
Add a dedicated discussion subsection titled something like:
\[
\text{``When EMPA reflects partition structure more strongly than perturbation magnitude.''}
\]

In that subsection:
\begin{itemize}
    \item acknowledge the phenomenon explicitly,
    \item explain whether it is due to the use of weight-only bias,
    \item test whether mean and covariance terms recover sensitivity,
    \item state whether the current score should be interpreted as a \emph{hybrid structural-discrepancy metric} rather than a pure budget signal.
\end{itemize}

\subsection*{Major Revision Item 24: Add audit-value correlation experiments}
\textbf{Problem.} The practical meaning of the audit score is still unclear.

\textbf{Required revision.}
Show correlation between audit output and actual privacy or utility outcomes, for example:
\begin{align}
\rho_1
&=
\mathrm{corr}
\bigl(
\mathrm{BAS}_\epsilon,
\mathrm{MIA\mbox{-}AUC}
\bigr), \\
\rho_2
&=
\mathrm{corr}
\bigl(
\mathrm{BAS}_\epsilon,
\mathrm{ReID\ Leakage}
\bigr), \\
\rho_3
&=
\mathrm{corr}
\bigl(
\mathrm{BAS}_\epsilon,
\Delta \mathrm{AP}
\bigr).
\end{align}

Even if these are only empirical proxies, they are much stronger than reporting feature-space discrepancies alone.

\section{VLM-Specific Section}

\subsection*{Major Revision Item 25: Reframe the VLM experiments as vision-tower auditing}
\textbf{Problem.} The current wording overstates the scope.

\textbf{Required revision.}
Throughout the VLM section, replace broad phrases such as ``auditing VLM privacy budgets'' with
\[
\text{auditing the visual encoder representations of VLMs}.
\]

Also add a formal scope statement:
\[
\text{Current study scope} = \text{vision tower only},
\qquad
\text{not full autoregressive multimodal generation}.
\]

\subsection*{Major Revision Item 26: Clarify that cross-method tables are scope comparisons, not head-to-head performance comparisons}
\textbf{Problem.} Comparing Bodhi against methods with different goals may be misleading.

\textbf{Required revision.}
Rename such tables using wording like:
\[
\text{``Scope and capability comparison with related privacy-aware vision/VLM methods.''}
\]
Do not present these as direct quantitative SOTA comparisons.

\section{Figures and Tables}

\subsection*{Major Revision Item 27: Redraw the main figures for readability}
\textbf{Problem.} The main pipeline figures appear overloaded.

\textbf{Required revision.}
Split the visual material into:
\begin{enumerate}
    \item one high-level framework overview,
    \item one BUA/TDA details figure,
    \item one EMPA fitting and discrepancy figure.
\end{enumerate}

\subsection*{Major Revision Item 28: Improve table interpretability}
\textbf{Required revision.}
For each quantitative table:
\begin{itemize}
    \item state the number of seeds,
    \item state whether lower or higher is better,
    \item explain normalization or scaling,
    \item highlight statistically significant best values only if significance is actually tested.
\end{itemize}

\section{Discussion and Limitations}

\subsection*{Major Revision Item 29: Expand the limitations section substantially}
\textbf{Problem.} The current limitations are present but not deep enough for a top-tier journal.

\textbf{Required revision.}
Add a structured limitations subsection including at least:

\paragraph{(a) Reference-model dependence.}
The audit depends on the evaluator-specified perturbation family:
\[
\mathcal{M}_{\mathrm{ref}}.
\]
A mismatch between actual and reference mechanisms can alter the interpretation of the score.

\paragraph{(b) Partition dependence.}
Audit quality depends on the quality of sensitive/non-sensitive grouping.

\paragraph{(c) Identifiability limits.}
Different perturbation settings may induce similar fitted mixture parameters.

\paragraph{(d) Scope limit for VLMs.}
Only the visual encoder is audited in the current work.

\paragraph{(e) No formal privacy guarantee.}
The framework is empirical and does not certify $(\epsilon,\delta)$-DP or equivalent formal notions.

\section{Conclusion}

\subsection*{Major Revision Item 30: Rewrite the conclusion to be precise and modest}
\textbf{Problem.} The conclusion should not restate broad claims that the body of the paper does not fully prove.

\textbf{Required revision.}
Conclude with a statement such as:
\begin{quote}
This work introduces an empirical feature-level auditing framework for assessing alignment between observed feature perturbations and evaluator-specified reference perturbation settings in vision backbones and VLM vision encoders.
\end{quote}

Then explicitly state the next steps:
\begin{itemize}
    \item stronger identifiability analysis,
    \item formal connection to privacy leakage,
    \item extension to full multimodal pipelines,
    \item calibration across heterogeneous perturbation families.
\end{itemize}

\section{New Equations and Definitions That Should Be Added}

\subsection*{1. Formal audit definition}
\begin{equation}
\mathcal{A}_{\text{audit}}
=
\mathcal{A}\!\left(
\tilde{f}, f, S_{\mathrm{sens}}, \epsilon ; \mathcal{M}_{\mathrm{ref}}, \Pi, \phi
\right).
\end{equation}

\subsection*{2. Reference-conditioned discrepancy target}
\begin{equation}
D_\epsilon
=
D\!\left(
\mathbb{P}(\tilde{f}\mid f,S_{\mathrm{sens}}),
\mathbb{P}(\tilde{f}^{\,\mathrm{ref}}_{\epsilon}\mid f,S_{\mathrm{sens}})
\right).
\end{equation}

\subsection*{3. Mixture-model representation of perturbation}
\begin{equation}
\mathbb{P}(\Delta f)
=
\sum_{k=1}^{K}
\lambda_k \,
\mathcal{N}(\Delta f; \mu_k, \Sigma_k),
\qquad
\Delta f = \tilde{f} - f.
\end{equation}

\subsection*{4. Full EMPA/BAS score}
\begin{equation}
\mathrm{BAS}_{\epsilon}
=
d_{\lambda}\!\left(
\lambda^{\mathrm{obs}}_{\epsilon},
\lambda^{\mathrm{ref}}_{\epsilon}
\right)
+
\alpha \,
d_{\mu}\!\left(
\mu^{\mathrm{obs}}_{\epsilon},
\mu^{\mathrm{ref}}_{\epsilon}
\right)
+
\beta \,
d_{\Sigma}\!\left(
\Sigma^{\mathrm{obs}}_{\epsilon},
\Sigma^{\mathrm{ref}}_{\epsilon}
\right).
\end{equation}

\subsection*{5. Rank-consistency metric for budget sensitivity}
\begin{equation}
\rho_{\mathrm{budget}}
=
\mathrm{corr}_{\mathrm{rank}}
\bigl(
\epsilon^{-1},
\mathrm{BAS}_{\epsilon}
\bigr).
\end{equation}

\subsection*{6. Correlation to downstream privacy/utility proxies}
\begin{equation}
\rho_{\mathrm{proxy}}
=
\mathrm{corr}
\bigl(
\mathrm{BAS}_{\epsilon},
\mathcal{R}_{\mathrm{privacy/utility}}
\bigr).
\end{equation}

\section{Priority Order for Revision}

\subsection*{Tier 1: Must fix before resubmission}
\begin{enumerate}
    \item Narrow the central claim across title, abstract, introduction, and conclusion.
    \item Formally define what is being audited.
    \item Add identifiability experiments for EMPA.
    \item Add task-relevant baselines.
    \item Add multi-seed statistics and significance testing.
\end{enumerate}

\subsection*{Tier 2: Strongly recommended}
\begin{enumerate}
    \item Validate BUA/TDA partition quality explicitly.
    \item Add mechanism-mismatch experiments.
    \item Add correlation analysis with actual privacy/utility proxies.
    \item Reframe VLM evaluation as vision-tower auditing only.
\end{enumerate}

\subsection*{Tier 3: Presentation improvements}
\begin{enumerate}
    \item Redraw figures.
    \item Simplify notation.
    \item Reorganize related work by task.
    \item Rename cross-method tables as scope comparisons.
\end{enumerate}

\section*{Bottom-Line Editorial Assessment}

In its current form, the paper presents an interesting and potentially useful auditing framework, but it does not yet meet top-tier journal standards because the main claim is broader than the evidence, the central EMPA quantity is not yet shown to be sufficiently identifiable, the experimental baselines are too weak, and the practical meaning of the audit score remains under-validated. The revision should therefore focus first on \textbf{problem-definition precision}, then on \textbf{identifiability and baseline strength}, and finally on \textbf{practical validation and statistical rigor}.